\section{袋外估计：只汇总没见过该样本的树}
\label{sec:tree-oob}
对$n$条记录有放回抽取$n$次构成一棵树的Bootstrap训练集。
某条记录可能被抽到多次，也可能一次都没抽到；后一种记录对这棵树称为袋外
（Out-of-Bag, OOB）。它不是对整个森林都未参与训练的独立测试记录。

\begin{figure}[htbp]
\centering\begin{tikzpicture}[x=13mm,y=8mm]
\foreach \x in {1,...,5}{\node[lab] at (\x,4.7) {样本\x};}
\foreach \y/\t in {4/1,3/2,2/3,1/4}{\node[lab,anchor=east] at (.3,\y) {树\t};}
\foreach \x/\y/\v in {1/4/2,2/4/1,4/4/2,2/3/1,3/3/2,5/3/2,1/2/1,2/2/2,3/2/1,5/2/1,1/1/1,3/1/1,4/1/2,5/1/1}{
 \node[draw=Teal,fill=Pale,minimum width=12mm,minimum height=7mm,font=\small] at (\x,\y) {\v};}
\foreach \x/\y in {3/4,5/4,1/3,4/3,4/2,2/1}{
 \node[draw=Amber,fill=Amber!15,minimum width=12mm,minimum height=7mm,font=\small] at (\x,\y) {0};}
\node[box,text width=43mm] at (8,2.5) {样本4的袋外预测\\仅来自树2、树3\\$(0.8+0.6)/2=0.7$};
\draw[Amber,thick] (4.42,1.55)--(4.5,1.55)--(4.5,3.45)--(4.42,3.45);
\draw[flow] (4.5,2.5)--(6.1,2.5);
\end{tikzpicture}

\caption{5条原记录、4份Bootstrap样本。格中为被抽次数；橙色0表示袋外。每行之和为5。}
\par\smallskip
\readfigure{对样本4，只有树2与树3没有用它训练。
若两树对该样本给正类概率0.8、0.6，袋外平均为0.7；树1、树4的预测不参与这一次袋外汇总。}
\end{figure}

每次未抽中某条记录的概率为$1-1/n$，独立抽$n$次后
\[
 P(\text{该记录袋外})=\left(1-\frac1n\right)^n\longrightarrow e^{-1}\approx0.368.
\]
$n=5$时是$(4/5)^5=0.32768$，不是精确36.8\%。
期望袋外记录数为$n(1-1/n)^n$；同一Bootstrap样本中不同记录的袋外指示并不独立。

\begin{enumerate}
\item 为每条记录找出没有抽到它的树，汇总这些树的预测。
\item 将袋外预测与该记录真实标签比较，再按指定指标汇总记录。
\item 检查袋外树数量；树太少时，部分记录可能没有袋外预测或估计很不稳定。
\end{enumerate}
不是先给每棵树算一个误差，再默认这些误差的简单均值就是同一个袋外指标。
分类平均概率还是硬投票，也需与实现约定一致。

\begin{pitfall}[袋外不自动评估整套预处理]
若特征筛选已利用全部训练标签，袋外记录的信息仍可能提前进入树。
袋外评分不能自动修复这种泄漏。评估完整流水线时，可采用折内拟合的交叉验证。
反复按袋外分数选参也会产生选择影响；分组和时间数据不能只靠逐行重采样模拟部署情形。
\end{pitfall}
